\documentclass{article}
\usepackage{amsthm}

\theoremstyle{plain}
\newtheorem{theorem}{Theorem}[section]
\newtheorem{proposition}[theorem]{Proposition}

\theoremstyle{definition}
\newtheorem{definition}[theorem]{Definition}
\newtheorem{example}[theorem]{Example}

\theoremstyle{remark}
\newtheorem{remark}[theorem]{Remark}

\begin{document}

\section{Various theorem styles}

\begin{definition}
\label{def:prime}
A prime is an integer greater than $1$ with no positive divisors other than $1$ and itself.
\end{definition}

\begin{proposition}
\label{prop:two-prime}
The number $2$ is prime.
\end{proposition}

\begin{proof}
Its only divisors are $1$ and $2$.
\end{proof}

\begin{remark}
\label{rem:two-only-even}
The number $2$ is the only even prime.
\end{remark}

\begin{example}
\label{ex:primes}
The first few primes are $2, 3, 5, 7, 11$.
\end{example}

\end{document}
